\documentclass[11pt]{article}

\usepackage[margin=1.1in]{geometry}
\usepackage{amsmath,amssymb,amsthm,mathtools}
\usepackage{booktabs}
\usepackage{enumitem}
\usepackage[colorlinks=true,linkcolor=blue,citecolor=blue,urlcolor=blue]{hyperref}

% ---------- theorem environments ----------
\theoremstyle{plain}
\newtheorem{theorem}{Theorem}[section]
\newtheorem{proposition}[theorem]{Proposition}
\newtheorem{lemma}[theorem]{Lemma}
\newtheorem{corollary}[theorem]{Corollary}
\newtheorem{fact}[theorem]{Fact} % imported results, stated without proof
\theoremstyle{definition}
\newtheorem{definition}[theorem]{Definition}
\newtheorem{assumption}[theorem]{Assumption}
\newtheorem{problem}[theorem]{Problem}
\theoremstyle{remark}
\newtheorem{remark}[theorem]{Remark}

% ---------- macros ----------
\newcommand{\E}{\mathbb{E}}
\newcommand{\Prob}{\mathbb{P}}
\newcommand{\R}{\mathbb{R}}
\newcommand{\N}{\mathbb{N}}
\newcommand{\one}{\mathbf{1}}
\newcommand{\cI}{\mathcal{I}}
\newcommand{\cL}{\mathcal{L}}
\newcommand{\cT}{\mathcal{T}}
\newcommand{\cS}{\mathcal{S}}
\newcommand{\cM}{\mathcal{M}}
\newcommand{\cH}{\mathcal{H}}
\newcommand{\cD}{\mathcal{D}}
\newcommand{\cA}{\mathcal{A}}
\newcommand{\cP}{\mathcal{P}}
\newcommand{\KL}{\mathrm{KL}}
\newcommand{\TV}{\mathrm{TV}}
\DeclareMathOperator*{\argmax}{arg\,max}

\title{A Mathematical Framework for Agent Families:\\
\large Library Learning as Budgeted Zeroth-Order Stochastic Optimization\\
under a Noisy, Adversarially-Maintained Instrument}
\author{Working notes --- agent-families project}
\date{June 2026 --- draft 0}

\begin{document}
\maketitle

\begin{abstract}
We formalize the agent-families system (DESIGN.md, Rev.~2) as a constrained
stochastic optimization problem over a discrete space of text libraries,
evaluated through a noisy and gameable instrument under a hard token budget.
Part~I states the system \emph{as it exists} in precise mathematical terms.
Part~II develops theory against that formalization: well-posedness of the
problem and of its behavioral objective (Banach fixed point); fundamental
sample-complexity limits on learning rate (Bretagnolle--Huber); anytime-valid
promotion gates with exact error control (Ville's inequality, e-value
composition, e-BH); unbiased causal credit (randomized masking) with
Monte-Carlo Shapley sample complexity; conditional approximation guarantees for
library growth (submodular greedy, with explicit noise degradation);
a convergence theorem for the learning loop (occupation-time bound on
non-stationary libraries); an exploitability bound for the grader as a
no-regret adversary; and the impossibility boundary (Rice, no-free-lunch,
Kolmogorov) that delimits what any such theory can claim. Every empirical
constant the theorems consume is mapped to a quantity the system already
measures (the Assumption Ledger, \S\ref{sec:ledger}).
\end{abstract}

\tableofcontents

%======================================================================
\section{Introduction, scope, and reading map}
%======================================================================

Three kinds of statement appear below, always labeled:
\begin{itemize}[leftmargin=2em]
  \item \textbf{Proved here.} Self-contained proofs from stated assumptions
        (most are short; the value is that the assumptions are \emph{ours},
        chosen to be enforceable or estimable by the harness).
  \item \textbf{Imported} (environment \emph{Fact}). Classical theorems used as
        black boxes, with citation.
  \item \textbf{Assumed} (environment \emph{Assumption}). Empirical hypotheses
        about LLM-dependent objects. These cannot be proved even in principle
        (\S\ref{sec:impossibility}); each is paired with its estimator in
        \S\ref{sec:ledger}.
\end{itemize}

The headline results:
\begin{center}
\begin{tabular}{@{}lll@{}}
\toprule
Claim & Result & Status \\
\midrule
Problem is well-posed & Prop.~\ref{prop:exists}, Thm.~\ref{thm:bisim} & proved \\
Learning rate is measurement-limited & Thm.~\ref{thm:lowerbound} & proved \\
Gates are valid under optional stopping & Thm.~\ref{thm:ville}, \ref{thm:ebh} & proved \\
Gates are sample-optimal & Fact~\ref{fact:wald}, Cor.~\ref{cor:orderopt} & imported + proved \\
Credit is unbiased and computable & Prop.~\ref{prop:ipw}, \ref{prop:shapley-mc} & proved \\
Library growth is near-optimal (conditional) & Thm.~\ref{thm:greedy}, Lem.~\ref{lem:epsgreedy} & proved, under Asm.~\ref{asm:submod} \\
The loop converges to statistical stationarity & Thm.~\ref{thm:drift} & proved \\
The grader's gameability is bounded & Thm.~\ref{thm:exploit} & proved \\
What cannot be proved & \S\ref{sec:impossibility} & impossibility \\
\bottomrule
\end{tabular}
\end{center}

%======================================================================
\part{The system as it exists}
%======================================================================

%----------------------------------------------------------------------
\section{Primitive spaces and the data model}\label{sec:primitives}
%----------------------------------------------------------------------

\begin{definition}[Insights and libraries]\label{def:library}
Fix a finite alphabet $\Sigma$ and a maximum insight length $\ell$. The
\emph{insight space} is $\cI = \Sigma^{\le \ell}$, a finite set. A
\emph{library state} is a pair $L = (I, \sigma)$ where $I \subseteq \cI$ is
finite and $\sigma : I \to \{\textsf{quarantined}, \textsf{active},
\textsf{dormant}, \textsf{retired}\}$ assigns lifecycle status. Write
$\mathrm{act}(L) = \sigma^{-1}(\textsf{active})$. The system additionally
maintains a partition of $\mathrm{act}(L)$ into skills and of skills into
agents; for the results below only the \emph{active multiset and its cap}
matter, so we suppress the partition except where noted. The \emph{library
space} is
\[
  \cL \;=\; \bigl\{\, L \;:\; |\mathrm{act}(L)| \le k \,\bigr\},
\]
where $k$ is the active cap (DESIGN \S4). Since $\cI$ is finite and statuses
are finitely many, $\cL$ is a \textbf{finite} set. This triviality does real
work in Proposition~\ref{prop:exists}.
\end{definition}

\begin{definition}[Retrieval]\label{def:retrieval}
A \emph{retrieval policy} is a map $R : Q \times \cL \to 2^{\cI}$ from queries
to subsets of active insights, subject to a token budget
$\sum_{i \in R(q,L)} |i| \le b$. In the implemented system $R$ is
embedding-KNN with an own-skills prior (DESIGN \S4); for the theory it is an
arbitrary measurable selection rule satisfying the budget.
\end{definition}

%----------------------------------------------------------------------
\section{The pipeline, episodes, and Ralph loops}\label{sec:pipeline}
%----------------------------------------------------------------------

\begin{definition}[Build pipeline as a stochastic kernel]\label{def:pipeline}
Let $\cT$ denote the set of \emph{targets} (apps presented through the
explorer's text channel) and $\cA$ the set of \emph{artifacts} (codebases).
The plan--work--verify pipeline with library $L$ defines a Markov kernel
\[
  A_L : \cT \to \Delta(\cA),
  \qquad
  A_L(T) = \mathrm{Law}\bigl(\text{artifact} \mid \text{target } T,\
  \text{library snapshot } L,\ \text{seed } \omega\bigr),
\]
where $\omega$ captures all sampling randomness (LLM decoding, scheduling).
The library is frozen within an episode (DESIGN \S11), so $A_L$ is a fixed
kernel per episode --- this is what makes the episode the unit of statistical
analysis.
\end{definition}

\begin{definition}[Ralph loop]\label{def:ralph}
Within one unit of work, iteration produces a sequence of states in a finite
set $\cS_{\mathrm{loop}} = \{\text{typed failure kinds}\} \cup
\{\textsf{pass}\}$, with $\textsf{pass}$ absorbing. We model the iteration as
a time-homogeneous Markov chain with transition matrix
$P = \begin{psmallmatrix} Q & r \\ 0 & 1 \end{psmallmatrix}$,
where $Q$ is substochastic over failure kinds
(Assumption~\ref{asm:markov}; Proposition~\ref{prop:absorb}).
\end{definition}

%----------------------------------------------------------------------
\section{The objective and the instrument}\label{sec:objective}
%----------------------------------------------------------------------

\begin{definition}[True grade and true objective]\label{def:objective}
Let $G^\ast : \cA \times \cT \to [0,1]$ denote the \emph{true grade}: the
behavioral-equivalence functional with its lexicographic quality tier
(formalized as a metric and a constrained program in
\S\ref{sec:wellposed}). The \emph{true objective} of a library is
\[
  J(L) \;=\; \E_{T \sim \cD}\;\E_{X \sim A_L(T)}\bigl[\,G^\ast(X, T)\,\bigr]
  \;\in\; [0,1],
\]
where $\cD$ is the (formally inaccessible) task distribution. $J$ is the
quantity the system is built to increase; it is \textbf{never observed}.
\end{definition}

\begin{definition}[Instrument]\label{def:instrument}
The grader provides, per episode on target $T$, an observation
\[
  \widehat{G}(X, T) \;=\; G^\ast(X, T) \;+\; b(X, T) \;+\; \varepsilon,
  \qquad \E[\varepsilon] = 0,\ \operatorname{Var}(\varepsilon) = \sigma^2,
\]
where $b$ is instrument bias (judge style bias, registry gaps; partially
removed by calibration, DESIGN \S10) and $\sigma^2$ decomposes by variance
components (judge $\times$ scenario $\times$ occasion $\times$
build-replicate; the G-theory model). All statistical guarantees below are
stated relative to this instrument model; \S\ref{sec:ledger} lists how $b$
and $\sigma^2$ are estimated by the frozen replay set, mutation audits, and
replicate episodes.
\end{definition}

%----------------------------------------------------------------------
\section{The learning operator and the optimization problem}\label{sec:problem}
%----------------------------------------------------------------------

\begin{definition}[Proposal--test--promote dynamics]\label{def:loop}
Learning proceeds in rounds $t = 0, 1, 2, \dots$ At round $t$:
\begin{enumerate}[label=(\roman*),leftmargin=2.5em]
  \item the reflector emits a candidate move
        $m_t \sim q(\cdot \mid \mathrm{trace}(L_t))$, where a \emph{move} is
        an admissible edit of $L_t$ (add a quarantined batch, retire, split,
        merge) and $q$ is an LLM-induced proposal distribution
        (Assumption~\ref{asm:proposal});
  \item a \emph{gate} evaluates the move by running episodes/benchmarks
        through the instrument and outputs $\textsf{accept}$ or
        $\textsf{revert}$;
  \item $L_{t+1} = L_t \oplus m_t$ if accepted, else $L_{t+1} = L_t$.
\end{enumerate}
This is exactly the quarantine $\to$ validate $\to$ promote lifecycle of
DESIGN \S5, viewed as a Markov chain on the finite space $\cL$.
\end{definition}

\begin{problem}[The system, as an optimization program]\label{prob:main}
\[
\begin{aligned}
  \text{maximize}_{\,\pi}\quad
    & \E\bigl[\,J(L_{\tau_B})\,\bigr] \\
  \text{subject to}\quad
    & L_{t+1} = L_t \oplus m_t \cdot \one\{\text{gate}_t = \textsf{accept}\},
      && m_t \sim q(\cdot \mid \mathrm{trace}(L_t)),\\
    & |\mathrm{act}(L_t)| \le k && \text{(cap)},\\
    & \textstyle\sum_{t \le \tau_B} c_t \le B && \text{(token budget)},
\end{aligned}
\]
where the policy $\pi$ chooses the gate design, the experiment allocation
(which targets, how many replications, when to stop each test), and the
governance moves; $c_t$ is the token cost of round $t$ and $\tau_B$ the
budget-exhaustion time. The objective $J$ is observable only through
$\widehat G$. \emph{No gradient of $J$ exists}: $\cL$ is discrete and $J$ has
no exploitable algebraic structure in $L$ beyond what is assumed in
\S\ref{sec:structure}.
\end{problem}

\begin{remark}
Problem~\ref{prob:main} identifies the system as \emph{budgeted zeroth-order
stochastic optimization with learned proposals over a finite combinatorial
space, evaluated by a biased noisy oracle}. Each clause picks out the body of
theory used in Part~II: finite space $\Rightarrow$ well-posedness
(\S\ref{sec:wellposed}); noisy oracle $\Rightarrow$ information bounds and
sequential tests (\S\ref{sec:measurement}); learned proposals $\Rightarrow$
drift analysis (\S\ref{sec:dynamics}); gameable oracle $\Rightarrow$ games and
regret (\S\ref{sec:adversarial}); set-valued decisions $\Rightarrow$
submodularity (\S\ref{sec:structure}).
\end{remark}

%======================================================================
\part{Theory developed against the formalization}
%======================================================================

%----------------------------------------------------------------------
\section{Well-posedness}\label{sec:wellposed}
%----------------------------------------------------------------------

\subsection{Existence}

\begin{proposition}[The problem is well-posed]\label{prop:exists}
$J : \cL \to [0,1]$ is well-defined and attains its maximum; an optimal
library $L^\ast \in \argmax_{L \in \cL} J(L)$ exists. Moreover every
quantity in Problem~\ref{prob:main} is measurable and bounded, so the
optimal value is finite and the expectations exist.
\end{proposition}

\begin{proof}
$G^\ast$ is bounded and measurable by construction
(Definition~\ref{def:objective}); $A_L(T)$ is a probability measure for each
$(L,T)$; hence the iterated expectation defining $J(L)$ exists and lies in
$[0,1]$. By Definition~\ref{def:library}, $\cL$ is finite; a real function on
a finite set attains its maximum.
\end{proof}

\begin{remark}
Uniqueness is neither available nor needed: the system seeks
\emph{improvement} under a budget, not the argmax. The substantive content of
well-posedness lives in the next subsection: that the behavioral part of
$G^\ast$ is itself a well-defined mathematical object.
\end{remark}

\subsection{The behavioral objective is a metric (and the suite estimates it)}

We model each app, at the granularity of the feature registry, as a finite
deterministic Mealy machine: states are abstracted UI/data states, inputs are
user actions, outputs are observations (rendered a11y-tree facts).

\begin{definition}[Behavioral distance operator]\label{def:bisimop}
Let $M_1 = (S_1, A, \delta_1, \lambda_1)$, $M_2 = (S_2, A, \delta_2,
\lambda_2)$ be finite Mealy machines over a common action alphabet $A$, with
outputs in a space $(O, d_O)$, $d_O \le 1$. Fix $\gamma \in (0,1)$. On the
set $\cP$ of pseudometrics $d : (S_1 \sqcup S_2)^2 \to [0,1]$, define
\[
  (F d)(s, t) \;=\; \max_{a \in A}\,
    \Bigl[\, (1-\gamma)\, d_O\bigl(\lambda(s,a), \lambda(t,a)\bigr)
      \;+\; \gamma\, d\bigl(\delta(s,a), \delta(t,a)\bigr) \,\Bigr].
\]
\end{definition}

\begin{theorem}[Existence, uniqueness, and meaning of the behavioral metric]
\label{thm:bisim}
$F$ has a unique fixed point $d_\gamma \in \cP$; iterates $F^n d_0$ converge
to it geometrically at rate $\gamma^n$ from any $d_0$. Moreover, for initial
states $s_0, t_0$:
\begin{enumerate}[label=(\alph*),leftmargin=2.5em]
  \item $d_\gamma(s_0, t_0) = 0$ if and only if the two machines are trace
        equivalent (identical outputs on every input word);
  \item if some input word $w = a_1 \cdots a_n$ produces outputs differing by
        $d_O$-distance $\eta$ at its last step, then
        $d_\gamma(s_0, t_0) \;\ge\; (1-\gamma)\,\gamma^{\,n-1}\eta$.
\end{enumerate}
\end{theorem}

\begin{proof}
\emph{Contraction.} For $d, d' \in \cP$ and any $(s,t)$, each bracketed term
differs between $d$ and $d'$ by at most $\gamma \|d - d'\|_\infty$, and
$\max_a$ is $1$-Lipschitz with respect to uniform perturbation of its
arguments; hence $\|F d - F d'\|_\infty \le \gamma \|d - d'\|_\infty$.
$\cP$ is a closed subset of the complete space of bounded functions on a
finite square under the sup norm (symmetry, the triangle inequality, and the
bound $[0,1]$ are preserved under uniform limits), and $F$ maps $\cP$ into
$\cP$: symmetry and boundedness are immediate; for the triangle inequality,
fix $s,u$ and let $a^\ast$ attain $(Fd)(s,u)$; then
\[
  (Fd)(s,u) = (1-\gamma) d_O(\lambda(s,a^\ast), \lambda(u,a^\ast))
            + \gamma\, d(\delta(s,a^\ast), \delta(u,a^\ast))
\]
is bounded, using the triangle inequalities of $d_O$ and $d$ through any
intermediate $t$, by the corresponding $a^\ast$-terms for $(s,t)$ and
$(t,u)$, hence by $(Fd)(s,t) + (Fd)(t,u)$. Banach's fixed point theorem gives
existence, uniqueness, and geometric convergence.

\emph{(a), $\Leftarrow$.} If $s_0, t_0$ are trace equivalent then so are
$\delta(s_0,a), \delta(t_0,a)$ for every $a$, with equal outputs. Starting
from $d_0 \equiv 0$, induction gives $(F^n d_0)(s,t) = 0$ for every
trace-equivalent pair and every $n$; the limit $d_\gamma$ vanishes there.

\emph{(a), $\Rightarrow$ and (b).} Unroll the fixed-point equation along
$w$: for each prefix, $d_\gamma(s,t) \ge (1-\gamma)\, d_O(\cdot) + \gamma\,
d_\gamma(\text{successors}) \ge \gamma\, d_\gamma(\text{successors})$, and at
the final step $d_\gamma \ge (1-\gamma)\eta$. Composing the $n-1$ discount
factors yields (b); in particular a distinguishing word forces
$d_\gamma(s_0,t_0) > 0$, which is the contrapositive of (a),
$\Rightarrow$.
\end{proof}

\begin{corollary}[What the grader's tiers estimate]\label{cor:estimators}
Each executed scenario is an input word; by Theorem~\ref{thm:bisim}(b) a
failed scenario yields a \emph{certified lower bound} on
$d_\gamma(\mathrm{target}, \mathrm{clone})$, while the registry suite as a
whole is a Monte-Carlo lower-bound estimator under the registry's sampling
measure. A falsifier (\S\ref{sec:adversarial}) estimates
$\sup_w$-type lower bounds. Quantitative (robustness-graded) verdicts refine
$\eta$ from $\{0,1\}$ to $[0,1]$, strictly increasing the information per
scenario at identical execution cost.
\end{corollary}

\begin{remark}[Modeling status]
Theorem~\ref{thm:bisim} is unconditional \emph{given} the Mealy abstraction;
the abstraction itself (finite registry-scoped state space, deterministic
transitions after seeding) is Assumption~\ref{asm:mealy}. Nondeterminism is
handled in practice by seeded resets; the probabilistic generalization
replaces $d$ on successors with a Wasserstein term and the same Banach
argument applies (imported: bisimulation metrics are optimal-transport
distances and are efficiently computable~\cite{ot-bisim}).
\end{remark}

%----------------------------------------------------------------------
\section{Measurement limits and valid gates}\label{sec:measurement}
%----------------------------------------------------------------------

\subsection{The learning rate is measurement-limited}

\begin{theorem}[Sample-complexity lower bound for detecting improvement]
\label{thm:lowerbound}
Let episode scores be observed through the instrument of
Definition~\ref{def:instrument} with noise dominated by a Gaussian of
variance $\sigma^2$ (conservative; subgaussian suffices with the same
constants). Any procedure that, from $n$ episodes, decides between
\[
  H_0 : J(L') \le J(L)
  \qquad\text{and}\qquad
  H_1 : J(L') \ge J(L) + \Delta
\]
with both error probabilities at most $\delta < 1/4$ requires
\[
  n \;\ge\; \frac{2\sigma^2}{\Delta^2}\,\ln\!\frac{1}{4\delta}.
\]
\end{theorem}

\begin{proof}
Reduce to the two simple hypotheses $P = \mathcal N(\mu, \sigma^2)^{\otimes n}$
and $Q = \mathcal N(\mu + \Delta, \sigma^2)^{\otimes n}$, for which
$\KL(P \,\|\, Q) = n\Delta^2 / (2\sigma^2)$. By the Bretagnolle--Huber
inequality~\cite{bh79,tsybakov}, every test $\psi$ satisfies
\[
  P(\psi = 1) + Q(\psi = 0) \;\ge\; \tfrac12\, e^{-\KL(P \| Q)}.
\]
If both errors are $\le \delta$ then $2\delta \ge \frac12
e^{-n\Delta^2/(2\sigma^2)}$, i.e.\
$n \ge \frac{2\sigma^2}{\Delta^2} \ln \frac{1}{4\delta}$.
\end{proof}

\begin{corollary}[Budgeted improvement bound]\label{cor:budget}
With per-episode cost $c$ and budget $B$, the number of improvement steps of
size $\Delta$ that can be \emph{verified} (at fixed $\delta$) during training
is at most
\[
  \frac{B}{c} \cdot \frac{\Delta^2}{2\sigma^2 \ln(1/4\delta)} .
\]
Halving instrument noise $\sigma$ is worth exactly a $4\times$ budget
multiplier. This is the formal sense in which the grader is the binding
component: \emph{no} change to the learner can beat this bound; only
instrument improvements move it.
\end{corollary}

\subsection{Anytime-valid gates}

\begin{definition}[e-process]\label{def:eprocess}
Let $\cH_0$ be a (possibly composite) null hypothesis. A nonnegative process
$(E_t)_{t \ge 0}$ adapted to the data filtration is an \emph{e-process} for
$\cH_0$ if $\E_P[E_\tau] \le 1$ for every $P \in \cH_0$ and every stopping
time $\tau$. (It suffices that $E_t$ is dominated by a nonnegative
supermartingale with initial value $1$ under every $P \in \cH_0$.)
\end{definition}

\begin{theorem}[Ville's inequality; validity of the promotion gate]
\label{thm:ville}
Let $(M_t)$ be a nonnegative supermartingale with $M_0 = 1$ under
$P \in \cH_0$. Then
\[
  \Prob_P\Bigl( \sup_{t \ge 0} M_t \ge 1/\alpha \Bigr) \;\le\; \alpha .
\]
Consequently the gate ``promote the batch the first time its e-process
reaches $1/\alpha$, with no constraint on when or whether you stop'' has
type-I error at most $\alpha$, \emph{uniformly over all stopping rules}.
\end{theorem}

\begin{proof}
Let $\tau = \inf\{t : M_t \ge 1/\alpha\}$. For each finite horizon $n$,
optional stopping for nonnegative supermartingales at the bounded time
$\tau \wedge n$ gives $\E[M_{\tau \wedge n}] \le \E[M_0] = 1$. On the event
$\{\tau \le n\}$ we have $M_{\tau \wedge n} = M_\tau \ge 1/\alpha$, so
$1 \ge \E[M_{\tau \wedge n}] \ge \alpha^{-1} \Prob(\tau \le n)$, i.e.\
$\Prob(\tau \le n) \le \alpha$ for every $n$. Monotone convergence as
$n \to \infty$ completes the proof.
\end{proof}

\begin{proposition}[Evidence composes across targets]\label{prop:eproduct}
If $E^{(1)}, \dots, E^{(m)}$ are e-values (e-processes at their stopping
times) for the same null, computed from \emph{independent} targets, then
$\prod_{j=1}^m E^{(j)}$ is an e-value. In particular ``the same insight
validated on $m$ independent targets'' multiplies evidence; the DESIGN
\S11 cross-target recurrence rule is the discretization of this statement.
\end{proposition}

\begin{proof}
By independence, $\E\bigl[\prod_j E^{(j)}\bigr] = \prod_j \E[E^{(j)}] \le 1$.
\end{proof}

\begin{theorem}[Batch promotion with false-discovery control: e-BH]
\label{thm:ebh}
Let $E_1, \dots, E_m$ be e-values for nulls $H_1, \dots, H_m$ (\emph{arbitrary
dependence} permitted). Order them decreasingly $E_{(1)} \ge \cdots \ge
E_{(m)}$, set $\hat k = \max\{k : E_{(k)} \ge \tfrac{m}{\alpha k}\}$ (and
$\hat k = 0$ if none), and reject the $\hat k$ hypotheses with the largest
e-values. Then $\mathrm{FDR} \le \alpha$.
\end{theorem}

\begin{proof}[Proof~\cite{ebh}]
Let $V$ be the number of true nulls rejected. If a true null $H_j$ is
rejected then $E_j \ge \frac{m}{\alpha \hat k}$, equivalently
$\frac{1}{\hat k} \le \frac{\alpha E_j}{m}$. Hence on $\{\hat k \ge 1\}$,
\[
  \frac{V}{\hat k}
  \;=\; \sum_{j \in \cH_0,\, \text{rejected}} \frac{1}{\hat k}
  \;\le\; \sum_{j \in \cH_0} \frac{\alpha E_j}{m},
\]
and taking expectations, $\mathrm{FDR} = \E[V / (\hat k \vee 1)] \le
\frac{\alpha}{m} \sum_{j \in \cH_0} \E[E_j] \le \alpha \frac{m_0}{m} \le
\alpha$.
\end{proof}

\begin{corollary}[False-promotion budget]\label{cor:falseprom}
If every promotion in rounds $1, \dots, T$ passes a level-$\alpha$ gate, the
expected number of promoted harmful batches is at most $\alpha T$ (linearity
of expectation; no independence needed). Under e-BH the realized
false-discovery \emph{rate} among promotions is $\le \alpha$ even with
arbitrarily dependent batches (parallel episodes against one snapshot).
\end{corollary}

\subsection{Efficiency of the gate}

\begin{fact}[Wald--Wolfowitz optimality of the SPRT~\cite{wald48}]
\label{fact:wald}
Among all tests (sequential or fixed-sample) of a simple null versus simple
alternative with type errors at most $(\alpha, \beta)$, the sequential
probability ratio test simultaneously minimizes $\E_0[N]$ and $\E_1[N]$, with
\[
  \E_1[N] \;\approx\; \frac{(1-\beta)\ln\frac{1-\beta}{\alpha} +
  \beta \ln\frac{\beta}{1-\alpha}}{\KL(P_1 \| P_0)} .
\]
\end{fact}

\begin{corollary}[The gate is order-optimal]\label{cor:orderopt}
For the Gaussian episode model, $\KL(P_1 \| P_0) = \Delta^2 / (2\sigma^2)$
per episode, so the SPRT/e-process gate decides in
$O\!\bigl(\tfrac{\sigma^2}{\Delta^2} \ln \tfrac1\alpha\bigr)$ expected
episodes --- matching the lower bound of Theorem~\ref{thm:lowerbound} up to
constants. \emph{No measurement policy can improve the exponent; only
$\sigma$ (instrument quality) and $\Delta$ (effect size per batch) are
levers.}
\end{corollary}

\subsection{Conformal triage of judge verdicts}

\begin{theorem}[Split-conformal coverage]\label{thm:conformal}
Let $S_1, \dots, S_n$ be nonconformity scores of calibration verdicts
(hand-adjudicated frozen replay pairs) and $S_{n+1}$ the score of a fresh
verdict, with $(S_1, \dots, S_{n+1})$ exchangeable. Let $\hat q$ be the
$\lceil (1-\alpha)(n+1) \rceil$-th smallest calibration score ($\hat q =
+\infty$ if the index exceeds $n$). Then
\[
  \Prob\bigl( S_{n+1} \le \hat q \bigr) \;\ge\; 1 - \alpha .
\]
Hence ``escalate to a panel exactly when $S_{n+1} > \hat q$'' guarantees that
non-escalated verdicts carry miscoverage at most $\alpha$, with no
distributional assumption on the judge.
\end{theorem}

\begin{proof}
By exchangeability the rank of $S_{n+1}$ among the $n+1$ scores (ties broken
uniformly) is uniform on $\{1, \dots, n+1\}$. The event $\{S_{n+1} \le \hat
q\}$ contains $\{\mathrm{rank}(S_{n+1}) \le \lceil (1-\alpha)(n+1) \rceil\}$,
whose probability is $\lceil (1-\alpha)(n+1)\rceil / (n+1) \ge 1-\alpha$.
\end{proof}

%----------------------------------------------------------------------
\section{Causal credit assignment}\label{sec:credit}
%----------------------------------------------------------------------

\begin{proposition}[Unbiased per-insight effects from randomized masking]
\label{prop:ipw}
In a quarantine validation run, mask each insight $i$ of the batch
independently with known probability $p \in (0,1)$ (include $i$ iff $Z_i =
1$). Let $Y_t(z)$ denote the (bounded) potential outcome of trial $t$ under
mask $z$, and define the estimand
$\tau_i = \E\bigl[ Y(Z_{-i}, 1) - Y(Z_{-i}, 0) \bigr]$
(the average effect of $i$ over the design distribution of its co-insights).
The Horvitz--Thompson estimator
\[
  \hat\tau_i \;=\; \frac1n \sum_{t=1}^n Y_t
  \Bigl( \frac{\one\{Z_{t,i}=1\}}{p} - \frac{\one\{Z_{t,i}=0\}}{1-p} \Bigr)
\]
satisfies $\E[\hat\tau_i] = \tau_i$ exactly, for every $n \ge 1$.
\end{proposition}

\begin{proof}
$Z_{t,i}$ is independent of $(Z_{t,-i}, Y_t(\cdot))$ by design. Hence
$\E\bigl[ Y_t \one\{Z_{t,i}=1\}/p \bigr]
 = \E\bigl[ Y_t(Z_{t,-i},1) \bigr] \E[\one\{Z_{t,i}=1\}]/p
 = \E[ Y_t(Z_{-i},1) ]$, and symmetrically for the second term. Subtract and
average.
\end{proof}

\begin{remark}
Observational fitness counters (wins/losses by co-occurrence) admit no such
statement: identification fails under co-retrieval confounding. Randomization
is the assumption, and it is one the \emph{scheduler} enforces, not one the
world must grant. This is the formal content of ``quarantine trials are
already an RCT in all but name.''
\end{remark}

\begin{fact}[Shapley uniqueness~\cite{shapley53}]\label{fact:shapley}
Let $v : 2^{[N]} \to \R$ be a coalition value (here: expected episode outcome
as a function of the active insight subset, estimated via
Proposition~\ref{prop:ipw}-style designs). The Shapley value is the unique
allocation satisfying efficiency, symmetry, dummy, and additivity.
\end{fact}

\begin{proposition}[Monte-Carlo Shapley sample complexity]\label{prop:shapley-mc}
Suppose $|v(S \cup \{i\}) - v(S)| \le r$ for all $S, i$. Sample $m$ uniform
permutations and average $i$'s marginal contributions. If
\[
  m \;\ge\; \frac{2 r^2}{\epsilon^2} \ln\frac{2N}{\delta},
\]
then with probability $\ge 1 - \delta$, simultaneously for all $N$ insights,
$|\hat\varphi_i - \varphi_i| \le \epsilon$.
\end{proposition}

\begin{proof}
Each permutation yields, for each $i$, an independent bounded draw with mean
$\varphi_i$ (the Shapley value is exactly the permutation-average of marginal
contributions). Hoeffding's inequality gives
$\Prob(|\hat\varphi_i - \varphi_i| > \epsilon) \le 2 e^{-m\epsilon^2 /
(2r^2)}$; a union bound over $N$ insights gives the claim.
\end{proof}

%----------------------------------------------------------------------
\section{Structural hypotheses: the library as submodular maximization}
\label{sec:structure}
%----------------------------------------------------------------------

The results in this section are \emph{conditional}: they trade one clearly
stated empirical hypothesis for strong guarantees. The hypothesis is
estimable from the system's own ablation data.

\begin{definition}[Submodularity; weak submodularity]
$f : 2^{\cI} \to \R_{\ge0}$ is \emph{monotone} if $f(S) \le f(T)$ for $S
\subseteq T$, and \emph{submodular} if for all $S \subseteq T$ and $e \notin
T$, $\;f(S \cup e) - f(S) \ge f(T \cup e) - f(T)$ (diminishing returns). For
$\gamma \in (0,1]$, $f$ is \emph{$\gamma$-weakly submodular} if for all
disjoint $S, U$:
$\;\sum_{e \in U} \bigl[f(S \cup e) - f(S)\bigr] \ge \gamma\,\bigl[f(S \cup
U) - f(S)\bigr]$.
\end{definition}

\begin{assumption}[Diminishing returns of insights]\label{asm:submod}
$J$, viewed as a function of the active insight set (other structure held
fixed), is monotone and $\gamma$-weakly submodular for some $\gamma > 0$.
\emph{Rationale}: redundant insights add less the fuller the library --- the
empirical phenomena the design already names (dedup pressure, shadowing,
cap necessity) are qualitative signatures of diminishing returns.
\emph{Estimator}: $\gamma$ is bounded from randomized-ablation data
(\S\ref{sec:ledger}).
\end{assumption}

\begin{theorem}[Greedy guarantee~\cite{nwf78}; proof included]\label{thm:greedy}
Let $f$ be monotone submodular with $f(\emptyset) = 0$ and let $S_k$ be the
greedy set after $k$ steps (each step adds the element of largest marginal
gain). Then for the optimum $S^\ast$ over sets of size $k$:
\[
  f(S_k) \;\ge\; \Bigl(1 - \tfrac1e\Bigr)\, f(S^\ast).
\]
Under $\gamma$-weak submodularity the bound becomes $1 - e^{-\gamma}$
\textup{(imported~\cite{daskempe})}.
\end{theorem}

\begin{proof}
Write $\delta_t = f(S^\ast) - f(S_t)$. By monotonicity and submodularity,
\[
  f(S^\ast) \;\le\; f(S_t \cup S^\ast)
  \;\le\; f(S_t) + \sum_{e \in S^\ast \setminus S_t}
          \bigl[f(S_t \cup e) - f(S_t)\bigr]
  \;\le\; f(S_t) + k \bigl[f(S_{t+1}) - f(S_t)\bigr],
\]
the last step because greedy picks the largest of at most $k$ such gains.
Rearranging, $\delta_{t+1} \le (1 - 1/k)\, \delta_t$, hence $\delta_k \le
(1-1/k)^k \delta_0 \le e^{-1} f(S^\ast)$.
\end{proof}

\begin{lemma}[Noise degrades greedy gracefully]\label{lem:epsgreedy}
If each greedy step selects an element whose true marginal gain is within
$\epsilon$ of the best (an ``$\epsilon$-approximate step''), then
\[
  f(S_k) \;\ge\; \bigl(1 - \tfrac1e\bigr) f(S^\ast) \;-\; k\,\epsilon .
\]
\end{lemma}

\begin{proof}
The recursion becomes $\delta_{t+1} \le (1 - 1/k)\delta_t + \epsilon$;
unrolling, $\delta_k \le e^{-1} f(S^\ast) + \epsilon \sum_{j=0}^{k-1} (1 -
1/k)^j \le e^{-1} f(S^\ast) + k\epsilon$.
\end{proof}

\begin{lemma}[Implementing $\epsilon$-approximate steps with the noisy oracle]
\label{lem:select}
Suppose marginal-gain evaluations are $\sigma$-subgaussian around their
means. Estimating each of $N$ candidates' gains with
\[
  n_0 \;\ge\; \frac{8\sigma^2}{\epsilon^2} \ln\frac{2N}{\delta'}
\]
episodes each, and selecting the empirical argmax, yields an
$\epsilon$-approximate step with probability $\ge 1 - \delta'$. Over $k$
steps, all succeed with probability $\ge 1 - k\delta'$
\textup{(union bound)}, so Lemma~\ref{lem:epsgreedy} applies with that
probability.
\end{lemma}

\begin{proof}
Hoeffding for subgaussian means: each estimate is within $\epsilon/2$ of its
mean with probability $\ge 1 - \delta'/N$; on the intersection (union bound)
the empirical argmax has true gain within $\epsilon$ of the best.
\end{proof}

\begin{fact}[The cap's tournament is streaming submodular maximization~\cite{badan14}]
\label{fact:streaming}
Insights arrive online; only $k$ may be retained. Threshold-based
swap rules achieve a $\tfrac12$-approximation to the best size-$k$ subset in
one pass with $O(k)$ memory, for monotone submodular objectives. Hence the
ratchet (tournament admission + displacement of the weakest incumbent), with
swap decisions made by the gates of \S\ref{sec:measurement}, inherits a
constant-factor guarantee under Assumption~\ref{asm:submod}.
\end{fact}

\begin{remark}[What this buys]
The cap stops being a defensive heuristic and becomes the memory parameter of
an online algorithm with a competitive ratio. The empirical task it induces
--- estimate $\gamma$, monitor whether marginal gains in fact diminish --- is
exactly the ablation telemetry the system already collects.
\end{remark}

%----------------------------------------------------------------------
\section{Dynamics: convergence of the learning loop}\label{sec:dynamics}
%----------------------------------------------------------------------

\subsection{The loop converges to statistical stationarity}

\begin{assumption}[Gate and proposal regularity]\label{asm:gate}
There exist $\Delta > 0$, $\alpha, \beta \in (0,1)$, $h \in (0,1]$ such that
for every round $t$ and every admissible move $m$:
\begin{enumerate}[label=(D\arabic*),leftmargin=3em]
  \item $J \in [0,1]$ (normalization).
  \item \emph{Level}: if $\Delta J(m \mid L_t) \le 0$ then
        $\Prob(\textsf{accept}) \le \alpha$
        \textup{(guaranteed by Theorem~\ref{thm:ville} given the instrument
        model)}. \emph{Power}: if $\Delta J(m \mid L_t) \ge \Delta$ then
        $\Prob(\textsf{accept}) \ge 1 - \beta$ \textup{(an experiment-design
        property: allocate episodes per Corollary~\ref{cor:orderopt})}.
  \item \emph{Bounded harm}: $\Delta J(m \mid L_t) \ge -h$ for every
        admissible move \textup{(enforced by bounded batch size, the joint
        confirmation run, and revertibility; $h$ is estimable from reverted
        batches)}.
\end{enumerate}
Define the \emph{proposal quality at scale $\Delta$}:
$\;p_\Delta(L) = \Prob_{m \sim q(\cdot \mid L)}\bigl( \Delta J(m \mid L) \ge
\Delta \bigr)$ --- the probability that the reflector can still find a
$\Delta$-improvement. This is the irreducibly empirical quantity
(Assumption~\ref{asm:proposal}).
\end{assumption}

\begin{theorem}[Occupation-time convergence]\label{thm:drift}
Under Assumption~\ref{asm:gate}, fix any $c > 0$ and define the
\emph{$(\Delta, c)$-stationary set}
\[
  \cS_c \;=\; \Bigl\{ L \in \cL \;:\;
  p_\Delta(L) \;\le\; \frac{(1+c)\,\alpha h}{(1-\beta)\,\Delta} \Bigr\}.
\]
Let $V_T = \#\{ t < T : L_t \notin \cS_c \}$ be the time spent at
non-stationary libraries. Then for every $T$:
\[
  \E[V_T] \;\le\; \frac{T}{c} \;+\; \frac{1 - J(L_0)}{c\,\alpha h},
  \qquad\text{hence}\qquad
  \limsup_{T \to \infty} \frac{\E[V_T]}{T} \;\le\; \frac1c .
\]
\end{theorem}

\begin{proof}
Decompose the one-step drift at $L_t = L$. Moves with $\Delta J \ge \Delta$
are proposed with probability $p_\Delta(L)$ and accepted with probability
$\ge 1 - \beta$, contributing at least $p_\Delta(L)(1-\beta)\Delta$. Moves
with $0 < \Delta J < \Delta$ contribute nonnegatively. Moves with $\Delta J
\le 0$ are accepted with probability $\le \alpha$ and each costs at most $h$,
contributing at least $-\alpha h$. Hence
\[
  \E\bigl[ J(L_{t+1}) - J(L_t) \mid L_t = L \bigr]
  \;\ge\; p_\Delta(L)\,(1-\beta)\,\Delta \;-\; \alpha h .
\]
If $L \notin \cS_c$, the right side is $\ge (1+c)\alpha h - \alpha h = c\,
\alpha h$; in all cases it is $\ge -\alpha h$. Telescoping over $t < T$ and
using $J \le 1$:
\[
  1 - J(L_0) \;\ge\; \E\bigl[J(L_T)\bigr] - J(L_0)
  \;\ge\; c\,\alpha h\, \E[V_T] \;-\; \alpha h \bigl( T - \E[V_T] \bigr)
  \;\ge\; c\,\alpha h\, \E[V_T] - \alpha h\, T .
\]
Rearranging yields the bound; divide by $T$ and let $T \to \infty$ for the
limit statement.
\end{proof}

\begin{remark}[Interpretation]
This is the zeroth-order, discrete analogue of ``SGD converges to first-order
stationary points'': global optimality is unavailable (and provably cannot be
certified, \S\ref{sec:impossibility}), but the chain provably spends all but
an arbitrarily small fraction of its time at libraries where
\emph{the reflector can no longer find improvements of size $\Delta$ except
with probability $O(\alpha h / ((1-\beta)\Delta))$}. Every knob is visible:
tightening $\alpha$ or boosting power $1-\beta$ shrinks the stationary set
(better final libraries) at a per-test episode cost given exactly by
Corollary~\ref{cor:orderopt}; larger detectable effects $\Delta$ make the
target easier but coarser. The theorem turns ``is it still learning?'' into
the measurable question ``has $p_\Delta(L_t)$ --- the validation pass rate of
fresh batches --- decayed to its floor?''
\end{remark}

\subsection{Inner-loop mathematics}

\begin{assumption}[Markov typed-failure dynamics]\label{asm:markov}
Within a Ralph loop, the typed verdict sequence is Markov with a
time-homogeneous transition matrix per (family, failure-kind) stratum.
\end{assumption}

\begin{proposition}[Expected iterations, in closed form]\label{prop:absorb}
Under Assumption~\ref{asm:markov}, write the chain as
$P = \begin{psmallmatrix} Q & r \\ 0 & 1 \end{psmallmatrix}$ with $Q$ the
substochastic transient block. If absorption is certain (spectral radius
$\rho(Q) < 1$), then $N = \sum_{j \ge 0} Q^j = (I - Q)^{-1}$ exists and the
expected iterations-to-pass from transient state $i$ is $(N \one)_i$;
absorption probabilities are $N r$. \textup{(Neumann series; standard, e.g.\
\cite{kemenysnell}.)} Consequently \textsc{max\_iters} and escalation
thresholds are functionals of an estimable matrix, not free constants.
\end{proposition}

\begin{fact}[Universal restarts are near-optimal~\cite{luby93}]\label{fact:luby}
For a Las Vegas process with unknown runtime distribution, the Luby schedule
achieves expected total time $O(\ell^\ast \log \ell^\ast)$, where $\ell^\ast$
is the expected time of the optimal fixed-cutoff restart policy, and no
universal schedule can do better than $\Omega(\ell^\ast \log \ell^\ast)$.
Fresh-context Ralph iterations approximately satisfy the i.i.d.\ restart
premise (Assumption ledger, \S\ref{sec:ledger}); the schedule then gives a
principled stuck-loop policy with a worst-case guarantee, replacing the flat
cap.
\end{fact}

%----------------------------------------------------------------------
\section{Adversarial grading: a quantitative anti-Goodhart bound}
\label{sec:adversarial}
%----------------------------------------------------------------------

\begin{definition}[The grading game]\label{def:game}
Fix a scenario/probe class $\cS$ (finite, or totally bounded with the bound
applied at covering resolution). Let $D(L, s) \in [0,1]$ denote the certified
divergence that probe $s$ extracts from the clone built with library $L$
(e.g.\ the robustness-violation magnitude, or the metric lower bound of
Corollary~\ref{cor:estimators}). At round $t$ the falsifier picks $s_t \in
\cS$ (knowing past observations), then observes $D(L_t, s_t)$. The
falsifier's \emph{external regret} is
\[
  R_T \;=\; \max_{s \in \cS} \sum_{t=1}^T D(L_t, s)
  \;-\; \sum_{t=1}^T D(L_t, s_t).
\]
\end{definition}

\begin{theorem}[Exploitability is bounded by the falsifier's regret]
\label{thm:exploit}
For any sequence of libraries $(L_t)$ --- in particular, one produced by an
optimizer adapting against the grader ---
\[
  \frac1T \sum_{t=1}^T D(L_t, s_t)
  \;\;\ge\;\;
  \max_{s \in \cS}\, \frac1T \sum_{t=1}^T D(L_t, s)
  \;-\; \frac{R_T}{T}.
\]
Hence if the \emph{measured} average divergence is kept below $\epsilon$, the
true average divergence against \emph{every fixed probe in $\cS$} is below
$\epsilon + R_T / T$. With a no-regret falsifier (e.g.\ EXP3 over $\cS$:
$R_T = O(\sqrt{T\, |\cS| \ln |\cS|})$~\cite{exp3}), the gameability of the
grader vanishes at rate $O(T^{-1/2})$.
\end{theorem}

\begin{proof}
Immediate rearrangement of Definition~\ref{def:game}.
\end{proof}

\begin{remark}[Sharpness of the caveat]
The guarantee is over the class $\cS$ the falsifier searches. A static test
suite is the degenerate case $|\cS| = $ suite, falsifier non-adaptive, $R_T =
\Theta(T)$: \emph{unbounded} exploitability --- the formal statement of ``a
fixed benchmark will be Goodharted.'' Growing $\cS$ by promoting discovered
counterexamples (the data flywheel) widens the quantifier of the theorem over
time. What no theorem can give is coverage of behaviors outside any probe
class ever searched; that residue is governed by
\S\ref{sec:impossibility}.
\end{remark}

%----------------------------------------------------------------------
\section{The impossibility boundary}\label{sec:impossibility}
%----------------------------------------------------------------------

These delimit, with standard results, the claims no extension of this
framework can make. They justify treating the corresponding objects as
\emph{estimated} rather than derived.

\begin{fact}[Rice's theorem $\Rightarrow$ exact grading is uncomputable]
Every nontrivial semantic property of programs is undecidable; in particular
exact behavioral equivalence of two programs is undecidable. Hence $G^\ast$
admits no exact decision procedure, and all grading is estimation of the
metric of Theorem~\ref{thm:bisim} relative to an observation measure (random
suite) or a search class (falsifier). ``Testing shows the presence, not the
absence, of bugs'' is the informal corollary.
\end{fact}

\begin{fact}[No free lunch $\Rightarrow$ structure must be assumed, then measured]
Averaged uniformly over all objectives, all black-box search algorithms have
identical performance~\cite{nfl}. Therefore \emph{any} claimed efficiency of
the loop (Theorems~\ref{thm:greedy}, \ref{thm:drift}) must rest on structural
hypotheses about $J$ --- here Assumptions~\ref{asm:submod} and
\ref{asm:proposal} --- and the scientifically honest posture is to estimate
those hypotheses from data rather than posit them silently. The Assumption
Ledger is the instrument of that honesty.
\end{fact}

\begin{fact}[Kolmogorov complexity is uncomputable $\Rightarrow$ MDL is code-relative]
No algorithm computes (or usefully lower-bounds, infinitely often) the
Kolmogorov complexity of strings. Hence ``the truly minimal library'' is not
a well-defined computational target; MDL-style governance is coherent only
relative to a fixed, declared coding scheme (Kraft inequality supplies the
accounting), and its role here is as a consistent \emph{criterion}, not an
optimality \emph{certificate}.
\end{fact}

\begin{assumption}[LLM-as-oracle]\label{asm:proposal}
The reflector's proposal distribution $q$, the judges' error rates, and the
resolver's fidelity have no derivable characterization. They enter every
theorem above only through scalar parameters
($p_\Delta(L)$, $\alpha$-validity of the e-construction, $\sigma$, $r$),
each of which the harness estimates (\S\ref{sec:ledger}). All guarantees are
conditional on those estimates in the usual sense of applied statistics.
\end{assumption}

\begin{assumption}[Mealy abstraction]\label{asm:mealy}
Registry-scoped app behavior is adequately modeled by a finite deterministic
Mealy machine after seeding/reset discipline. Violations (true
nondeterminism, hidden state) degrade Theorem~\ref{thm:bisim} to its
probabilistic generalization~\cite{ot-bisim}.
\end{assumption}

%----------------------------------------------------------------------
\section{The Assumption Ledger}\label{sec:ledger}
%----------------------------------------------------------------------

Every empirical constant consumed by Part II, with the telemetry that
estimates it. This table is the contract between the theory and the running
system: each row is a number the harness must produce before the
corresponding guarantee is invoked outside shadow mode.

\begin{center}
\small
\begin{tabular}{@{}p{0.30\textwidth}p{0.22\textwidth}p{0.40\textwidth}@{}}
\toprule
\textbf{Assumed quantity} & \textbf{Used in} & \textbf{Estimator in the system} \\
\midrule
Instrument noise $\sigma^2$ (by variance component) &
Thm.~\ref{thm:lowerbound}, Lem.~\ref{lem:select}, Cor.~\ref{cor:orderopt} &
Replicate episodes at fixed snapshot; G-theory fit on frozen replay set \\
\addlinespace
Judge error rates / e-construction validity &
Thm.~\ref{thm:ville}, \ref{thm:ebh} &
Mutation-seeded audits (false-pass rate); frozen replay re-judging \\
\addlinespace
Exchangeability of judge scores &
Thm.~\ref{thm:conformal} &
Drift monitors on the calibration stream; re-split on prompt-set change \\
\addlinespace
Bounded harm $h$ per accepted move &
Thm.~\ref{thm:drift} &
Outcome distribution of reverted batches; joint-confirmation failures \\
\addlinespace
Power $1-\beta$ at effect $\Delta$ &
Thm.~\ref{thm:drift} &
Designed from $\hat\sigma$ via Cor.~\ref{cor:orderopt} (episode allocation) \\
\addlinespace
Proposal quality $p_\Delta(L)$ &
Thm.~\ref{thm:drift} (stationarity diagnosis) &
Validation pass rate of fresh reflector batches per epoch \\
\addlinespace
Weak-submodularity constant $\gamma$ &
Thm.~\ref{thm:greedy}, Fact~\ref{fact:streaming} &
Randomized-ablation marginal-gain curvature (Prop.~\ref{prop:ipw} data) \\
\addlinespace
Markov property of typed failures &
Prop.~\ref{prop:absorb} &
Order-1 vs.\ order-2 likelihood-ratio test on span sequences \\
\addlinespace
I.i.d.\ restart premise &
Fact~\ref{fact:luby} &
Correlation of successive fresh-context iteration outcomes \\
\addlinespace
Probe-class coverage $\cS$ &
Thm.~\ref{thm:exploit} &
Irreducible; mitigated by counterexample promotion (growing $\cS$) \\
\bottomrule
\end{tabular}
\end{center}

%----------------------------------------------------------------------
\section{Outlook: what a paper can claim}\label{sec:outlook}
%----------------------------------------------------------------------

The defensible composite claim, assembled from the parts:

\begin{quote}
\itshape
The agent-families loop is a well-posed budgeted zeroth-order optimization
(Prop.~\ref{prop:exists}) of a well-defined behavioral objective
(Thm.~\ref{thm:bisim}). Its measurement layer is provably valid under
continuous monitoring and optional stopping (Thms.~\ref{thm:ville},
\ref{thm:ebh}, \ref{thm:conformal}) and order-optimal in sample complexity
(Thm.~\ref{thm:lowerbound} with Cor.~\ref{cor:orderopt}). Its credit
assignment is exactly unbiased by design (Prop.~\ref{prop:ipw}) with
controlled approximation cost (Prop.~\ref{prop:shapley-mc}). Conditional on a
measured diminishing-returns constant, its governance is
constant-factor-competitive (Thm.~\ref{thm:greedy},
Fact~\ref{fact:streaming}). The loop provably concentrates its time on
libraries at which no detectable improvement remains proposable
(Thm.~\ref{thm:drift}), and the grader's gameability decays at the
falsifier's regret rate (Thm.~\ref{thm:exploit}). The residue --- proposal
quality, oracle fidelity, out-of-class behavior, and global optimality ---
is provably beyond proof (\S\ref{sec:impossibility}) and is therefore
estimated, continuously, by the Assumption Ledger.
\end{quote}

Open problems worth working: (i) estimate $\gamma$ and test
Assumption~\ref{asm:submod} on real ablation data --- the single hypothesis
carrying the most theory; (ii) sharpen Theorem~\ref{thm:drift} with
state-dependent harm under the rollback mechanism (the $-h$ bound is
loose when reversion is fast); (iii) extend Theorem~\ref{thm:exploit} to
bandit feedback over a structured continuous probe space via covering
numbers; (iv) a two-timescale stochastic-approximation treatment of
calibrating $\widehat G$ while $L_t$ moves (Borkar conditions); (v) the
restless-bandit (Whittle) formulation of curriculum as a corollary layer over
Theorem~\ref{thm:lowerbound}'s per-target information accounting.

%======================================================================
\begin{thebibliography}{99}\small

\bibitem{bh79} J.~Bretagnolle and C.~Huber.
Estimation des densit\'es: risque minimax.
\emph{Z.\ Wahrscheinlichkeitstheorie verw.\ Gebiete}, 47:119--137, 1979.

\bibitem{tsybakov} A.~Tsybakov.
\emph{Introduction to Nonparametric Estimation}. Springer, 2009. (Ch.~2 for
Bretagnolle--Huber and two-point lower bounds.)

\bibitem{ville} J.~Ville.
\emph{\'Etude critique de la notion de collectif}. Gauthier-Villars, 1939.
(Ville's inequality.)

\bibitem{ramdas} A.~Ramdas, P.~Gr\"unwald, V.~Vovk, G.~Shafer.
Game-theoretic statistics and safe anytime-valid inference.
\emph{Statistical Science}, 38(4), 2023. (Modern e-process treatment.)

\bibitem{ebh} R.~Wang and A.~Ramdas.
False discovery rate control with e-values.
\emph{J.\ Royal Statistical Society B}, 84(3):822--852, 2022.

\bibitem{wald48} A.~Wald and J.~Wolfowitz.
Optimum character of the sequential probability ratio test.
\emph{Ann.\ Math.\ Statist.}, 19(3):326--339, 1948.

\bibitem{vovk-conformal} V.~Vovk, A.~Gammerman, G.~Shafer.
\emph{Algorithmic Learning in a Random World}. Springer, 2005. (Conformal
prediction.)

\bibitem{neyman} J.~Neyman.
On the application of probability theory to agricultural experiments (1923).
Translated in \emph{Statistical Science}, 5(4):465--472, 1990.

\bibitem{shapley53} L.~S.~Shapley.
A value for $n$-person games. In \emph{Contributions to the Theory of Games
II}, 307--317. Princeton, 1953.

\bibitem{nwf78} G.~Nemhauser, L.~Wolsey, M.~Fisher.
An analysis of approximations for maximizing submodular set functions---I.
\emph{Mathematical Programming}, 14:265--294, 1978.

\bibitem{daskempe} A.~Das and D.~Kempe.
Submodular meets spectral: greedy algorithms for subset selection, sparse
approximation and dictionary selection. \emph{ICML}, 2011. (Weak
submodularity / submodularity ratio.)

\bibitem{badan14} A.~Badanidiyuru, B.~Mirzasoleiman, A.~Karbasi, A.~Krause.
Streaming submodular maximization: massive data summarization on the fly.
\emph{KDD}, 2014.

\bibitem{kemenysnell} J.~Kemeny and J.~Snell.
\emph{Finite Markov Chains}. Springer, 1976. (Absorbing-chain fundamental
matrix.)

\bibitem{luby93} M.~Luby, A.~Sinclair, D.~Zuckerman.
Optimal speedup of Las Vegas algorithms.
\emph{Information Processing Letters}, 47(4):173--180, 1993.

\bibitem{exp3} P.~Auer, N.~Cesa-Bianchi, Y.~Freund, R.~Schapire.
The nonstochastic multiarmed bandit problem.
\emph{SIAM J.\ Computing}, 32(1):48--77, 2002.

\bibitem{ot-bisim} S.~Calo et al.
Bisimulation metrics are optimal transport distances, and can be computed
efficiently. \emph{NeurIPS}, 2024. arXiv:2406.04056.

\bibitem{nfl} D.~Wolpert and W.~Macready.
No free lunch theorems for optimization.
\emph{IEEE Trans.\ Evolutionary Computation}, 1(1):67--82, 1997.

\bibitem{golovin} D.~Golovin and A.~Krause.
Adaptive submodularity: theory and applications in active learning and
stochastic optimization. \emph{JAIR}, 42:427--486, 2011. (For the outlook's
VOI-greedy claim.)

\bibitem{borkar} V.~Borkar.
\emph{Stochastic Approximation: A Dynamical Systems Viewpoint}. Cambridge,
2008. (Two-timescale SA, for the outlook.)

\end{thebibliography}

\end{document}
